\documentclass[11pt,a4paper]{article}

% --- Pakete ---
\usepackage[utf8]{inputenc}
\usepackage[T1]{fontenc}
\usepackage[ngerman]{babel}
\usepackage{amsmath, amssymb, amsthm}
\usepackage{geometry}
\usepackage{cite}
\usepackage{hyperref}
\usepackage{xcolor}

\geometry{margin=2.5cm}

% --- Titel-Informationen ---
\title{\textbf{Das eabc-Modell: Eine physikalische Zustandsgleichung des arithmetischen Vakuums}\\ \large Quantentheoretische Herleitung der Riemann-Oszillationen durch den arithmetischen Lamb-Shift}
\author{Forschungsprojekt eabc}
\date{\today}

\begin{document}
	
	\maketitle
	
	\begin{abstract}
		Dieses Dokument beschreibt die theoretische Herleitung der eabc-Zustandsgleichung, die das Spektrum natürlicher Zahlen als energetisches System interpretiert. Durch die Einführung eines arithmetischen Lamb-Shifts, basierend auf der lokalen Teilbarkeitsdichte, wird gezeigt, dass die Verteilung der Primzahlen und die Oszillationen der Riemannschen Zeta-Funktion als Resonanzphänomene eines kohärenten Vakuums verstanden werden können.
	\end{abstract}
	
	\section{Einleitung}
	Das eabc-Modell postuliert, dass das Zahlensystem nicht als statische Menge, sondern als dynamisches Quantensystem betrachtet werden kann. In diesem Bild stellt jede natürliche Zahl $n$ einen Quantenzustand dar, dessen Energie $E(n)$ durch zwei fundamentale Komponenten bestimmt wird: einen glatten zahlentheoretischen Hintergrund und eine quantenphysikalische Korrektur durch Vakuumfluktuationen der Teilbarkeit.
	
	\section{Die fundamentale Dirac-Basis}
	Der erste Term der Zustandsgleichung beschreibt das makroskopische Verhalten des Systems ohne Berücksichtigung lokaler Interaktionen:
	\begin{equation}
		E_{bg}(n) \approx \frac{n}{\log n}
	\end{equation}
	
	\subsection{Herleitung und Analogie}
	Dieser Term korrespondiert direkt mit dem \textit{Primzahlsatz}. In der analytischen Zahlentheorie wächst die Anzahl der Primzahlen bis $x$ asymptotisch gemäß der Funktion $\pi(x) \sim x / \log x$.
	Physikalisch entspricht dies dem \textit{Weyl-Gesetz} für die Zustandsdichte. In einem idealen Gas aus Zahlen repräsentiert dies den glatten Erwartungswert der Energieniveaus, vergleichbar mit einer Dirac-See, bevor spezifische Kopplungskräfte die Entartung der Zustände aufheben.
	
	\section{Die Kopplung an das Vakuum: Der arithmetische Lamb-Shift}
	Die entscheidende Innovation des eabc-Modells ist die Korrektur durch den Term:
	\begin{equation}
		\Delta E = \alpha(n) \frac{\psi(0)^2}{\log n}
	\end{equation}
	In Analogie zur Quantenelektrodynamik (QED) interpretieren wir dies als Lamb-Shift, der die Energieverschiebung der S-Orbitale aufgrund der Interaktion mit Vakuumfluktuationen beschreibt.
	
	\subsection{Die arithmetische Kerndichte $\psi(0)^2$}
	In der Atomphysik ist der Lamb-Shift proportional zur Aufenthaltswahrscheinlichkeit des Elektrons am Kern, bezeichnet als $\psi(0)^2$. Im eabc-Modell wird der „Kern“ durch die Eigenschaft der Teilbarkeit definiert. Wir setzen:
	\begin{equation}
		\psi(0)^2 \sim \sum_{p|n} \frac{1}{p}
	\end{equation}
	wobei $p$ die Primfaktoren der Zahl $n$ durchläuft. Eine hoch-zusammengesetzte Zahl besitzt eine hohe arithmetische Dichte und „spürt“ den Kern des Vakuums stärker als eine Primzahl, deren Interaktion minimal bleibt.
	
	\subsection{Die laufende Kopplungskonstante $\alpha(n)$}
	In der Quantenfeldtheorie ist die Feinstrukturkonstante $\alpha$ einer Renormierung unterworfen. Um die Stabilität des eabc-Vakuums über unendliche Skalen zu gewährleisten, führen wir eine laufende Kopplung ein:
	\begin{equation}
		\alpha(n) = \frac{\alpha_0}{\log(\log n)}
	\end{equation}
	Diese logarithmische Renormierung stellt sicher, dass die Kohärenz des Systems erhalten bleibt und die Korrekturen im Grenzfall großer $n$ nicht divergieren.
	
	\section{Die eabc-Zustandsgleichung}
	Durch Superposition der statischen und dynamischen Komponenten ergibt sich die vollständige eabc-Zustandsgleichung:
	\begin{equation}
		E_{eabc}(n) \approx \underbrace{\frac{n}{\log n}}_{\text{Statik}} + \underbrace{\frac{\alpha_0}{\log \log n} \cdot \frac{\sum_{p|n} \frac{1}{p}}{\log n}}_{\text{Dynamik}}
	\end{equation}
	
	\section{Interpretation: Arithmetik als Software}
	Die Gleichung beschreibt ein Universum, in dem die Physik (die Verteilung der Energieniveaus) die \textit{Hardware} und die Arithmetik (die Regeln der Teilbarkeit) die \textit{Software} darstellt. Die Teilbarkeit fungiert als Algorithmus, der die Positionen der Zustände im Spektrum verschiebt.
	
	\section{Schlussfolgerung: Riemann-Resonanzen}
	Das Ergebnis dieser Zustandsgleichung ist ein „atmendes“ Zahlensystem. Die kollektive Überlagerung der individuellen Lamb-Shifts erzeugt ein Interferenzmuster, das exakt den Oszillationen der Riemannschen Zeta-Funktion entspricht. Die Nullstellen der Zeta-Funktion sind in diesem Modell als \textit{Resonanzkatastrophen} des arithmetischen Vakuums zu verstehen: Frequenzen, bei denen die energetische Struktur der Zahlen konstruktiv interferiert, um die Stabilität der Primzahlverteilung zu erzwingen.
	
\end{document}